% ==============================================================
% Projektive Strukturtheorie (PST)
% Dokument PST-3-P: Projektionstheorie
% Version 0.1
% Kompiliert mit: pdfLaTeX
% ==============================================================

\documentclass[11pt,a4paper]{article}

% ---------- Grundpakete ----------
\usepackage[utf8]{inputenc}
\usepackage[T1]{fontenc}
\usepackage[german]{babel}
\usepackage{amsmath,amssymb,amsthm}
\usepackage{geometry}
\usepackage{parskip}
\usepackage{enumitem}
\usepackage{booktabs}
\usepackage{xcolor}
\usepackage{hyperref}

% ---------- Layout ----------
\geometry{margin=2.5cm}
\setlist{itemsep=0.2em, topsep=0.2em}
\hypersetup{
    colorlinks=true,
    linkcolor=blue,
    urlcolor=blue,
    pdftitle={PST-3-P: Projektionstheorie -- Version 0.1},
    pdfauthor={Forschungsprogramm Ontoph\"{a}nomenalismus}
}

% ---------- Theorem-Umgebungen ----------
\theoremstyle{plain}
\newtheorem{theorem}{Theorem}[section]
\newtheorem{proposition}{Proposition}[section]
\newtheorem{lemma}{Lemma}[section]
\newtheorem{corollary}{Korollar}[section]

\theoremstyle{definition}
\newtheorem{definition}{Definition}[section]
\newtheorem{remark}{Bemerkung}[section]
\newtheorem{methodennote}{Methodennote}[section]
\newtheorem{hypothese}{Arbeitshypothese}[section]
\newtheorem{prinzip}{Prinzip}[section]

% Neue Kommandos
\newcommand{\OPP}{\mathcal{M}_{\mathrm{OPP}}}
\newcommand{\Kmat}{\mathbf{K}}
\newcommand{\deff}{d_{\mathrm{eff}}}
\newcommand{\Proj}{\mathcal{P}}

% ---------- Titel ----------
\title{
  \textbf{Projektive Strukturtheorie (PST)}\\[0.4em]
  \textbf{Dokument PST-3-P: Projektionstheorie}\\[0.4em]
  {\large Eigenschaften physikalischer Projektionen, Symmetriebrechung,
  Dimensionsemergenz und der universelle Operator $\mathcal{P}^\ast$}\\[0.5em]
  {\normalsize Version~0.1}
}
\author{
  \textbf{Forschungsprogramm Ontoph\"{a}nomenalismus}\\
  \small Siehe TRANSPARENCY f\"{u}r methodologische Urheberschaft
  und KI-Einsatz.
}
\date{Juli 2026}

\begin{document}

\maketitle

\begin{abstract}
Dieses Dokument entwickelt die Projektionstheorie der PST. Es beantwortet
die zentrale Frage: \emph{Welche Eigenschaften muss ein Projektionsoperator
besitzen, damit er physikalisch sinnvolle Strukturen erzeugt?}

Die drei Kernaussagen des Dokuments sind:

\begin{enumerate}
  \item Nicht jede Projektion ist physikalisch --- physikalische
        Projektionen sind eine ausgezeichnete Klasse.
  \item Projektion ist Symmetriebrechung: Der OPP besitzt maximale
        Symmetrie; die physikalische Raumzeit entsteht durch
        Einschr\"{a}nkung dieser Symmetrie.
  \item $\deff \approx 4$ ist kein Zufall --- es gibt strukturelle
        Gr\"{u}nde warum genau vier Dimensionen emergieren.
\end{enumerate}

Die Suche nach dem universellen Projektionsoperator $\Proj^\ast$ ist
die zentrale offene Frage der PST.

Es gelten die methodologischen Konventionen gem\"{a}\ss{} MKO.
\end{abstract}

\tableofcontents
\newpage

%==============================================================
\section{Die zentrale Frage: Welche Projektionen sind physikalisch?}
%==============================================================

\begin{methodennote}[Paradigmenwechsel in der PST]
Die fr\"{u}he Phase der PST fragte: \glqq{}Hat der OPP vier
Dimensionen?\grqq{} Die Antwort war: Nein, der rohe OPP besitzt
$\deff \approx 30$--$45$.

Die reife Phase fragt: \glqq{}Welche Projektionen erzeugen
vierdimensionale Physik?\grqq{} Das ist eine fundamentale Umkehrung:
Nicht der OPP \emph{ist} vierdimensional --- die Physik
\emph{erscheint} vierdimensional weil der Projektionsoperator
$\Proj$ genau die Strukturen hervorhebt, die wir als Raumzeit
wahrnehmen.
\end{methodennote}

Die Frage nach dem physikalischen Projektionsoperator ist nicht
\"{a}quivalent zur Frage nach dem OPP. Es sind zwei verschiedene
Ebenen:

\begin{center}
\begin{tabular}{lll}
  \toprule
  \textbf{Ebene} & \textbf{Objekt} & \textbf{Frage} \\
  \midrule
  Ontisch & OPP & Was existiert? \\
  Projektiv & $\Proj$ & Wie wird es sichtbar? \\
  Ph\"{a}nomenal & Raumzeit & Was beobachten wir? \\
  \bottomrule
\end{tabular}
\end{center}

%==============================================================
\section{Algebraische Eigenschaften physikalischer Projektionen}
%==============================================================

\begin{definition}[Projektionsoperator]
Ein \textbf{Projektionsoperator} $\Proj\colon \OPP \to
\mathcal{M}_{\mathcal{O}}$ bildet den hochdimensionalen OPP auf
einen niedrigdimensionalen Beobachtungsraum $\mathcal{M}_{\mathcal{O}}$
ab. Die algebraischen Eigenschaften (vgl.\ DPR):
\begin{enumerate}
  \item \textbf{Idempotenz:} $\Proj \circ \Proj = \Proj$ --- eine
        bereits projizierte Struktur \"{a}ndert sich durch erneute
        Projektion nicht.
  \item \textbf{Nicht-Injektivit\"{a}t:} $\Proj(x) = \Proj(y)
        \not\Rightarrow x = y$ --- verschiedene OPP-Zust\"{a}nde
        k\"{o}nnen auf denselben Beobachtungszustand abgebildet werden.
  \item \textbf{Koh\"{a}renzerhaltung:} Relational verbundene Zust\"{a}nde
        im OPP erscheinen auch im Beobachtungsraum als verbunden.
\end{enumerate}
\end{definition}

\begin{definition}[Physikalische Projektion]
Eine Projektion $\Proj$ hei\ss{}t \textbf{physikalisch}, wenn sie
zus\"{a}tzlich folgende Eigenschaften erf\"{u}llt:
\begin{enumerate}
  \item \textbf{Globale Kohärenzerhaltung:} $\Proj$ erh\"{a}lt globale
        Strukturen des OPP, nicht nur lokale Details.
  \item \textbf{Stabilit\"{a}t:} $\deff(\Proj)$ ist stabil gegen\"{u}ber
        kleinen St\"{o}rungen des Projektionsoperators.
  \item \textbf{Lorentz-Invarianz (Hypothese):} Die emergierte Metrik
        besitzt Lorentz-Symmetrie --- $(+,-,-,-)$ Signatur.
  \item \textbf{Eich-Kompatibilit\"{a}t (Hypothese):} Die Eichgruppen
        des Standardmodells $SU(3) \times SU(2) \times U(1)$ bleiben
        als Symmetriegruppen der Projektion erhalten.
\end{enumerate}
\end{definition}

\begin{remark}[Warum globale Kohärenz entscheidend ist]
Numerische Experimente (PST-4-G) zeigen klar: Projektionsverfahren
die \emph{lokale} Strukturen bevorzugen (PCA, Random Projection) liefern
$\deff < 4$ oder instabile Ergebnisse. Verfahren die
\emph{globale} Strukturen erhalten (Diffusion Maps, Kernel PCA) liefern
$\deff \approx 4$.

Dies suggeriert: Physikalische Projektionen sind genau jene, die die
globale Relationsstruktur des OPP kohär\"{a}rent auf den Beobachtungsraum
\"{u}bertragen.
\end{remark}

%==============================================================
\section{Projektion als Symmetriebrechung}
%==============================================================

\begin{prinzip}[Projektion = Symmetriebrechung]
Der OPP besitzt eine maximale Symmetriegruppe $G_{\OPP}$ ---
die Gruppe aller Transformationen die die Relationsstruktur invariant
lassen. Die physikalische Raumzeit entsteht durch Einschr\"{a}nkung
dieser Symmetrie:
\[
  G_{\OPP} \;\xrightarrow{\Proj}\; G_\Proj
  = \Proj(G_{\OPP}) \;\subseteq\; G_{\OPP}.
\]
\end{prinzip}

\begin{remark}[Lorentz-Gruppe als projizierte Symmetrie]
Die Lorentz-Gruppe $SO(1,3)$ ist in dieser Sichtweise keine
fundamentale Symmetrie des Universums, sondern eine
\emph{projizierte Symmetrie} --- die Symmetriegruppe die nach der
Anwendung von $\Proj^\ast$ auf den OPP verbleibt.

Das ist eine tiefgreifende Umkehrung: Nicht der OPP muss
Lorentz-Symmetrie besitzen --- die Projektion muss sie erzeugen.
\end{remark}

\begin{hypothese}[Eichgruppen als Projektivsymmetrien --- MKO Ebene~5]
Die Eichgruppen des Standardmodells
$SU(3) \times SU(2) \times U(1)$ k\"{o}nnten Symmetriegruppen sein,
die nach der physikalischen Projektion im Beobachtungsraum verbleiben:
\[
  \Proj^\ast(G_{\OPP}) \supseteq SU(3) \times SU(2) \times U(1)
  \times SO(1,3).
\]
Falls dies zutrifft, erkl\"{a}rt die PST nicht nur die Raumzeit, sondern
auch die Materiestruktur des Universums als projektive Erscheinung.

Diese Hypothese ist spekulativ (MKO Ebene~5) und erfordert erhebliche
weitere mathematische Arbeit.
\end{hypothese}

\begin{remark}[Verbindung zu Kaluza-Klein]
Es gibt eine strukturelle \"{A}hnlichkeit zur Kaluza-Klein-Theorie: Dort
werden Eichsymmetrien als geometrische Symmetrien zus\"{a}tzlicher
kompakter Raumdimensionen gedeutet. In der PST wird der umgekehrte
Weg gegangen: Aus dem hochdimensionalen OPP werden durch Projektion
sowohl die 3+1 Raumzeit als auch die Eichstrukturen herausdestilliert.
Ob dies mathematisch \"{a}quivalent ist oder eine neue Perspektive
er\"{o}ffnet, bleibt eine Frage der PST-6-R.
\end{remark}

%==============================================================
\section{Warum vier Dimensionen?}
%==============================================================

\begin{methodennote}[Offenste Frage der PST]
Die Frage warum genau $\deff \approx 4$ emergiert --- und nicht 3, 5
oder 11 --- ist die tiefste offene Frage der PST. Numerisch ist das
Ergebnis reproduzierbar. Warum es so ist, ist noch nicht verstanden.
\end{methodennote}

\begin{hypothese}[Stabilitätshypothese f\"{u}r $d = 4$ --- MKO Ebene~5]
Genau vier Dimensionen sind die einzige Zahl f\"{u}r die die folgende
Kombination von Eigenschaften gleichzeitig m\"{o}glich ist:
\begin{enumerate}
  \item \textbf{Kausale Struktur:} $d = 4$ ist die minimale Dimension
        f\"{u}r eine nicht-triviale kausale Struktur (Lichtkegelgeometrie).
  \item \textbf{Lorentz-Invarianz:} $SO(1,3)$ besitzt genau in $d = 4$
        bestimmte Eigenschaften (Weyl-Spinoren, selbstduale Formen), die
        f\"{u}r das Standardmodell unabdingbar sind.
  \item \textbf{Stabilit\"{a}t von Bahnen:} Planetenbahnen sind nur in
        $d = 3$ Raumdimensionen stabil. Mehr Dimensionen f\"{u}hren zu
        Instabilit\"{a}ten.
  \item \textbf{Eichtheorie-Kompatibilit\"{a}t:} Yang-Mills-Theorien sind
        in $d = 4$ renormierbar.
\end{enumerate}
Falls der OPP durch Diffusion Maps projiziert wird, k\"{o}nnte das
Minimum der Freien Energie der Projektion bei $\deff = 4$ liegen ---
als Folge der globalen Struktur des OPP.
\end{hypothese}

\begin{remark}[Anthropischer Einwand und seine Grenze]
Man k\"{o}nnte einwenden: Die vier Dimensionen werden beobachtet weil
wir in einem Universum mit vier Dimensionen leben (anthropisches
Argument). Die PST versucht, dieses Argument zu \"{u}berwinden: Nicht
\glqq{}Wir beobachten 4D weil wir 4D-Wesen sind\grqq{}, sondern
\glqq{}Der universelle Projektionsoperator $\Proj^\ast$ erzeugt
notwendigerweise 4D, unabh\"{a}ngig vom Beobachter.\grqq{} Das ist eine
st\"{a}rkere und pr\"{u}fbarere These.
\end{remark}

%==============================================================
\section{Der universelle Projektionsoperator $\Proj^\ast$}
%==============================================================

\begin{definition}[Universeller Projektionsoperator]
Der \textbf{universelle Projektionsoperator} $\Proj^\ast$ ist der
(noch zu findende) Operator der:
\begin{enumerate}
  \item Notwendigerweise $\deff = 4$ erzeugt (nicht als Annahme,
        sondern als Eigenschaft des OPP).
  \item Die Lorentz-Signatur $(+,-,-,-)$ reproduziert.
  \item Die Eichgruppen $SU(3) \times SU(2) \times U(1)$ erzeugt.
  \item Quantisierung als Konsequenz seiner endlichen Aufl\"{o}sung liefert.
  \item Die bekannten Naturkonstanten reproduziert.
\end{enumerate}
\end{definition}

\begin{remark}[Status von $\Proj^\ast$]
$\Proj^\ast$ ist das Fernziel der PST. Kein existierender Formalismus
erf\"{u}llt alle f\"{u}nf Bedingungen. Diffusion Maps bei $t \approx 9$--$10$
erf\"{u}llt Bedingung~1 numerisch --- die restlichen vier Bedingungen
sind noch offen.
\end{remark}

\begin{prinzip}[Randbedingungsprinzip]
Die Suche nach $\Proj^\ast$ folgt dem Top-Down-Ansatz: Die bekannte
Physik (Standardmodell, Lorentz-Invarianz, Quantenmechanik) liefert
die Randbedingungen. Ein Projektionsoperator der diese Randbedingungen
erf\"{u}llt, ist ein Kandidat f\"{u}r $\Proj^\ast$. Die PST sucht den
\emph{einfachsten} Operator der alle Randbedingungen erf\"{u}llt.
\end{prinzip}

%==============================================================
\section{Quantisierung als projektives Artefakt}
%==============================================================

\begin{hypothese}[Quantisierung aus endlicher Projektionsaufl\"{o}sung
--- MKO Ebene~5]
Quantisierung ist keine ontologische Eigenschaft des OPP, sondern
eine Eigenschaft der Projektion: Weil der Projektionsoperator
$\Proj$ eine endliche Informationsaufl\"{o}sung besitzt, entstehen
diskrete Zust\"{a}nde im projizierten Raum.

Formal: Die minimale Informationseinheit $I_0$ (PST-2-M) entspricht
dem Planck'schen Wirkungsquantum $\hbar$:
\[
  I_0 \;\leftrightarrow\; \hbar.
\]
Unterhalb von $I_0$ gibt es keine projektive Bestimmtheit --- und damit
keine unterscheidbaren Quantenzust\"{a}nde. Die Quantisierung von Energie,
Impuls und Drehimpuls w\"{a}re dann eine direkte Konsequenz der minimalen
Projektionsaufl\"{o}sung.

Diese Hypothese ist spekulativ (MKO Ebene~5) und erfordert eine
pr\"{a}zise mathematische Ausarbeitung in PST-6-R.
\end{hypothese}

\begin{remark}[Verbindung zur Unsch\"{a}rferelation]
Die Nicht-Injektivit\"{a}t des Projektionsoperators ($\Proj(x) = \Proj(y)
\not\Rightarrow x = y$, vgl.\ DPR) k\"{o}nnte eine konzeptuelle Wurzel
der quantenmechanischen Unsch\"{a}rfe sein: Verschiedene OPP-Zust\"{a}nde
sind projektiv nicht unterscheidbar. Das entspricht der Heisenbergschen
Unsch\"{a}rferelation: Orts- und Impulsmessung k\"{o}nnen nicht gleichzeitig
beliebig pr\"{a}zise sein --- weil der Projektionsoperator beide nicht
gleichzeitig vollst\"{a}ndig aufl\"{o}st.
\end{remark}

%==============================================================
\section{Lorentz-Signatur als projektive Eigenschaft}
%==============================================================

\begin{hypothese}[Lorentz-Signatur aus Projektion --- MKO Ebene~5]
Die charakteristische Minkowski-Signatur $(+,-,-,-)$ der Raumzeit ist
keine Eigenschaft des OPP, sondern des Projektionsoperators $\Proj^\ast$.

Konkret: Die Zeitdimension und die drei Raumdimensionen entstehen durch
unterschiedliche Rollen im Projektionsspektrum --- der erste Eigenwert
$\lambda_1$ der Gram-Matrix k\"{o}nnte die zeitartige Richtung
definieren, w\"{a}hrend $\lambda_2, \lambda_3, \lambda_4$ die raumartigen
Richtungen definieren.

Die negativen Vorzeichen in der Lorentz-Signatur k\"{o}nnten aus dem
Imagin\"{a}rteil von Eigenvektoren emergieren --- bei komplexen
Kernel-Matrizen. Dies ist eine offene Forschungshypothese f\"{u}r PST-6-R.
\end{hypothese}

%==============================================================
\section{Beobachter als Projektionsoperatoren}
%==============================================================

\begin{proposition}[Beobachter = Projektionsoperatoren]
Im Rahmen der PST ist ein physikalischer Beobachter identisch mit
einem Projektionsoperator $\Proj_{\mathcal{O}}$:
\[
  \text{Beobachter } \mathcal{O}
  \;\equiv\;
  \Proj_{\mathcal{O}}\colon \OPP \to \mathcal{M}_{\mathcal{O}}.
\]
Verschiedene Beobachter entsprechen verschiedenen Projektionsoperatoren
und \glqq{}sehen\grqq{} deshalb verschiedene (aber konsistente)
Ausschnitte des OPP.
\end{proposition}

\begin{remark}[Verbindung zu DPR und RdE]
Dies ist die mathematische Ausarbeitung der \emph{Relativit\"{a}t der
Existenz} (RdE) aus DPR: Raumzeitliche Existenz ist relativ zum
Projektionsoperator. Was f\"{u}r $\Proj_{\mathcal{O}_1}$ existiert, muss
nicht f\"{u}r $\Proj_{\mathcal{O}_2}$ existieren --- wegen der
Nicht-Injektivit\"{a}t.
\end{remark}

\begin{remark}[Verbindung zu DFD]
Das Fragedifferential $\mathcal{D}(E) = D_{\mathrm{KL}}(Q_E \| P)$
misst die Informationsdistanz zwischen dem projizierten Zustand $Q_E$
eines Beobachters und dem Gesamtfeld $P$. In der PST-Sprache: Es misst
den \glqq{}Abstand\grqq{} zwischen dem Projektionsoperator
$\Proj_{\mathcal{O}}$ und dem universellen Operator $\Proj^\ast$.
\end{remark}

%==============================================================
\section{Die Hierarchie der Projektioneigenschaften}
%==============================================================

Nicht alle Projektioneigenschaften sind gleich fundamental. Wir
unterscheiden:

\begin{enumerate}
  \item \textbf{Notwendige Eigenschaften:} Idempotenz, Nicht-Injektivit\"{a}t,
        globale Koh\"{a}renzerhaltung. Jede physikalische Projektion
        muss diese besitzen.
  \item \textbf{Hinreichende Eigenschaften (Hypothese):} Lorentz-Invarianz,
        Eich-Kompatibilit\"{a}t, $\deff = 4$. Diese charakterisieren
        $\Proj^\ast$ spezifisch.
  \item \textbf{Emergente Eigenschaften:} Quantisierung, Kausalit\"{a}t,
        Dynamik. Diese folgen aus den notwendigen und hinreichenden
        Eigenschaften --- sie werden nicht separat postuliert.
\end{enumerate}

%==============================================================
\section{Offene Fragen der Projektionstheorie}
%==============================================================

\begin{enumerate}
  \item \textbf{$\Proj^\ast$ aus Prinzipien:} Kann der universelle
        Projektionsoperator aus dem Prinzip der minimalen Ontologie
        und der Randbedingung Physik eindeutig hergeleitet werden?
  \item \textbf{Lorentz-Signatur:} Wie emergiert die
        $(+,-,-,-)$-Signatur aus einer signaturfreien Gram-Matrix?
  \item \textbf{Eichgruppen:} Warum $SU(3) \times SU(2) \times U(1)$
        und nicht andere Gruppen?
  \item \textbf{Quantisierung:} Ist $I_0 \leftrightarrow \hbar$
        mathematisch pr\"{a}zisierbar?
  \item \textbf{Eindeutigkeit:} Gibt es genau einen $\Proj^\ast$ oder
        eine Familie von Operatoren die alle physikalisch \"{a}quivalent
        sind?
  \item \textbf{Dynamik:} Wie emergiert zeitliche Evolution aus dem
        statischen OPP? Ist die Schr\"{o}dinger-Gleichung eine
        projektive Eigenschaft?
\end{enumerate}

%==============================================================
\section{Schnittstellen}
%==============================================================

\begin{itemize}
  \item \textbf{PST-1-F:} Das Rekonstruktionsziel OPP $\rightarrow
        \mathcal{P} \rightarrow$ Raumzeit $\rightarrow$ Physik wird
        hier mathematisch ausgearbeitet.
  \item \textbf{PST-2-M:} Die Gram-Matrix und Diffusion Maps aus
        PST-2-M sind die mathematischen Grundlagen der Projektionstheorie.
  \item \textbf{DPR (Die Projektive Reduktion):} PST-3-P ist die
        mathematische Ausarbeitung von DPR. Idempotenz, Nicht-Injektivit\"{a}t
        und RdE aus DPR werden hier formalisiert.
  \item \textbf{DFD (Das Fragedifferential):} Das Fragedifferential
        misst in PST-Sprache den Abstand vom universellen Operator
        $\Proj^\ast$.
  \item \textbf{PST-4-G:} Die numerischen Evidenzen f\"{u}r die
        Projektionstheorie werden in PST-4-G dokumentiert.
  \item \textbf{PST-6-R:} Die vollst\"{a}ndige physikalische Ausarbeitung
        von Lorentz-Signatur, Eichgruppen und Quantisierung obliegt
        PST-6-R.
\end{itemize}

\end{document}

